\documentclass[12pt]{article}
\usepackage[utf8]{inputenc}
\usepackage{amsmath, amssymb, geometry, graphicx, setspace, xcolor}
\geometry{letterpaper, margin=1in}
\setstretch{1.15}

\title{\textbf{Instructor Solution Manual: \\ Computational Astrodynamics Lab (NASA GMAT)}}
\author{Module: Space Engineering Vehicle and Mission Design}
\date{}

\begin{document}

\maketitle

\section*{Instructor Notes \& Constants}
To ensure the analytical hand-calculations match the GMAT Differential Corrector (DC) outputs perfectly, students must use the exact constants GMAT uses for a Point Mass Earth. 
\begin{itemize}
    \item Earth Gravitational Parameter: $\mu_{\oplus} = 398600.4415 \text{ km}^3/\text{s}^2$
    \item Earth Equatorial Radius: $R_{\oplus} = 6378.136 \text{ km}$
\end{itemize}
\textit{Grading Tip:} GMAT's DC should converge to within $10^{-4}$ km/s of these analytical results if the force model was properly simplified (Point Mass only, no drag).

\vspace{0.5cm}

% ---------------------------------------------------------
% SOLUTION 1: HOHMANN TRANSFER
% ---------------------------------------------------------
\section*{Solution 1: LEO to GEO Hohmann Transfer}

\textbf{Given Parameters:}
\begin{itemize}
    \item $r_1 = 6378.136 + 300 = 6678.136 \text{ km}$
    \item $r_2 = 6378.136 + 35786 = 42164.136 \text{ km}$
\end{itemize}

\subsection*{Part A: Analytical Solution}
\begin{enumerate}
    \item \textbf{Semi-major axis of transfer ellipse ($a_{tx}$):}
    \begin{equation*}
        a_{tx} = \frac{r_1 + r_2}{2} = \frac{6678.136 + 42164.136}{2} = \textbf{24421.136 \text{ km}}
    \end{equation*}

    \item \textbf{Initial parking orbit velocity ($v_1$):}
    \begin{equation*}
        v_1 = \sqrt{\frac{\mu_{\oplus}}{r_1}} = \sqrt{\frac{398600.4415}{6678.136}} = \textbf{7.7258 \text{ km/s}}
    \end{equation*}

    \item \textbf{Periapsis velocity of the transfer orbit ($v_{tx,1}$):}
    \begin{equation*}
        v_{tx,1} = \sqrt{\mu_{\oplus} \left( \frac{2}{r_1} - \frac{1}{a_{tx}} \right)} = \sqrt{398600.4415 \left( \frac{2}{6678.136} - \frac{1}{24421.136} \right)} = \textbf{10.1487 \text{ km/s}}
    \end{equation*}

    \item \textbf{First Impulsive Burn ($\Delta V_1$):}
    \begin{equation*}
        \Delta V_1 = v_{tx,1} - v_1 = 10.1487 - 7.7258 = \textbf{2.4229 \text{ km/s}}
    \end{equation*}

    \item \textbf{Apoapsis velocity of the transfer orbit ($v_{tx,2}$):}
    \begin{equation*}
        v_{tx,2} = \sqrt{\mu_{\oplus} \left( \frac{2}{r_2} - \frac{1}{a_{tx}} \right)} = \sqrt{398600.4415 \left( \frac{2}{42164.136} - \frac{1}{24421.136} \right)} = \textbf{1.6074 \text{ km/s}}
    \end{equation*}

    \item \textbf{Final circular GEO velocity ($v_2$):}
    \begin{equation*}
        v_2 = \sqrt{\frac{\mu_{\oplus}}{r_2}} = \sqrt{\frac{398600.4415}{42164.136}} = \textbf{3.0746 \text{ km/s}}
    \end{equation*}

    \item \textbf{Second Impulsive Burn ($\Delta V_2$):}
    \begin{equation*}
        \Delta V_2 = v_2 - v_{tx,2} = 3.0746 - 1.6074 = \textbf{1.4672 \text{ km/s}}
    \end{equation*}

    \item \textbf{Total Mission $\Delta V$:}
    \begin{equation*}
        \Delta V_{total} = 2.4229 + 1.4672 = \textbf{3.8901 \text{ km/s}}
    \end{equation*}
\end{enumerate}

\textcolor{blue}{\textbf{Instructor Check for Part B (GMAT):}} Students should report $\Delta V_1 \approx 2.4229$ km/s and $\Delta V_2 \approx 1.4672$ km/s from the GMAT message window. If their numbers differ by more than a few meters per second, they likely forgot to change the Earth gravity model from JGM-2 (spherical harmonics) to Point Mass.

\newpage

% ---------------------------------------------------------
% SOLUTION 2: PURE PLANE CHANGE
% ---------------------------------------------------------
\section*{Solution 2: Pure Orbital Plane Change}

\textbf{Given Parameters:}
\begin{itemize}
    \item $r = 6378.136 + 1200 = 7578.136 \text{ km}$
    \item $\Delta i = 28.5^\circ - 0^\circ = 28.5^\circ$
\end{itemize}

\subsection*{Part A: Analytical Solution}
\begin{enumerate}
    \item \textbf{Orbital velocity ($v_c$):}
    \begin{equation*}
        v_c = \sqrt{\frac{\mu_{\oplus}}{r}} = \sqrt{\frac{398600.4415}{7578.136}} = \textbf{7.2526 \text{ km/s}}
    \end{equation*}

    \item \textbf{Required Impulsive Maneuver ($\Delta V$):}
    \begin{equation*}
        \Delta V = 2 v_c \sin \left( \frac{\Delta i}{2} \right) = 2 (7.2526) \sin \left( \frac{28.5^\circ}{2} \right) 
    \end{equation*}
    \begin{equation*}
        \Delta V = 14.5052 \times \sin(14.25^\circ) = 14.5052 \times 0.24615 = \textbf{3.5705 \text{ km/s}}
    \end{equation*}
\end{enumerate}

\subsection*{Part B: GMAT Verification}
\textcolor{blue}{\textbf{Instructor Check:}} The student must answer the conceptual question about the sign of the normal burn. At the \textbf{descending node}, the spacecraft is crossing the equator from North to South. To flatten the orbit to the equator (reduce inclination), the $\Delta V$ must be applied in the \textbf{Positive Normal} direction (pointing "up" or North). Therefore, the GMAT output for `Element2` must be approximately $\mathbf{+3.5705 \text{ km/s}}$.

\newpage

% ---------------------------------------------------------
% SOLUTION 3: PATCHED CONIC DEPARTURE
% ---------------------------------------------------------
\section*{Solution 3: Interplanetary Departure (Patched Conic)}

\textbf{Given Parameters:}
\begin{itemize}
    \item Target $v_\infty = 3.5 \text{ km/s}$
    \item $r_p = 6378.136 + 400 = 6778.136 \text{ km}$
\end{itemize}

\subsection*{Part A: Analytical Solution}
\begin{enumerate}
    \item \textbf{Specific Energy ($C_3$):}
    \begin{equation*}
        C_3 = v_\infty^2 = 3.5^2 = \textbf{12.25} 
    \end{equation*}

    \item \textbf{Parking orbit velocity ($v_c$):}
    \begin{equation*}
        v_c = \sqrt{\frac{\mu_{\oplus}}{r_p}} = \sqrt{\frac{398600.4415}{6778.136}} = \textbf{7.6686 \text{ km/s}}
    \end{equation*}

    \item \textbf{Required burnout velocity at periapsis ($v_p$):}
    \begin{equation*}
        v_p = \sqrt{v_\infty^2 + \frac{2\mu_{\oplus}}{r_p}} = \sqrt{12.25 + \frac{2(398600.4415)}{6778.136}}
    \end{equation*}
    \begin{equation*}
        v_p = \sqrt{12.25 + 117.6141} = \sqrt{129.8641} = \textbf{11.3958 \text{ km/s}}
    \end{equation*}

    \item \textbf{Required prograde engine burn ($\Delta V_{depart}$):}
    \begin{equation*}
        \Delta V_{depart} = v_p - v_c = 11.3958 - 7.6686 = \textbf{3.7272 \text{ km/s}}
    \end{equation*}
\end{enumerate}

\subsection*{Part B: GMAT Verification}
\textcolor{blue}{\textbf{Instructor Check:}} The DC should converge to $\Delta V \approx 3.7272 \text{ km/s}$. 
\textit{Common Student Error:} If the student propagated to the SOI using \texttt{Sat.ElapsedDays = 10} instead of using the radial stopping condition \texttt{Sat.Earth.RMAG = 925000}, their $C_3$ target might not converge properly because the spacecraft hasn't reached the asymptotic limit. Ensure their GMAT script utilizes the precise RMAG stopping condition.

\end{document}